%% notes-tex/figure-3-gate/sections/6-overnight.tex
%% SOURCE: tables/v22_arms.tex, tables/v22_runs.tex and tables/shrinkage.tex are generated by
%% experiments/019-simb-multimodal/scripts/v22_tables.py from results/v22_readout.json
%% (v22_readout.py, W&B project torchcell_019_expr_v22) and
%% results/prediction_shrinkage_probe.json (prediction_shrinkage_probe.py). The timing table
%% is hand-authored from the Lightning profiler tables and stack samples recorded in
%% notes/experiments.019-simb-multimodal.scripts.gh_profile_cgt_expr.md (GilaHyper jobs 3302
%% to 3327) and from W&B perf/epoch_seconds. Delta figures are from the job logs and W&B
%% project torchcell_019_expr_v21, read 2026-10-06 between 02:50 and 10:30 CDT.

\clearpage
\section{The night of 2026-10-05: speed, the Delta round, and the fast round v22\secstatus{todo}}
\label{sec:night}

\pfstep{Correction: v21 on Delta cannot reach its window as launched} \Cref{sec:status}
projected 22 to 24 hours per pack from the training progress bar, 49 to 59 s. The logged
epoch time is 177 s on every run: the reference arm was at epoch 69 after 3.3 hours of
training. At that pace 1{,}200 epochs take 59 hours against the 48 hour limit, and a pack
stops near epoch 945, short of the registered window of epochs 1{,}000 to 1{,}200. The cause
was read from forty stack dumps of one run over six minutes: in all forty the main thread
sat in the eval-mode train pass, spawning a data-loader worker. That pass rebuilt its loader
every tenth epoch, and on Delta a spawned worker re-imports the software stack from Lustre,
more than twelve minutes each time. It follows from launching the round with one loader
worker; with zero the pass would have been slow but would not have respawned anything.
\src{W\&B torchcell_019_expr_v21, perf/epoch_seconds; Delta job 22690625}

\pfstep{Where an epoch went} A profile of the reference run on the GilaHyper cards, then
stack samples of live runs, located five costs, none of them the model's own arithmetic
(\cref{tab:night-timing}). Producing a sample cost 50 ms (a database read, a JSON parse, the
reconstruction of an experiment record with thousands of phenotype values) and was repeated
for every sample every epoch. Decoding each strain's target row from the batch ran as a
Python loop over the strains with about ten points per strain where the CPU waits on the
GPU, 59 to 65 percent of a training process's time once several processes shared a card.
The perturbation operator ran one strain at a time. The epoch metrics ranked a 1{,}100 by
6{,}127 matrix on the CPU three times an epoch. And all six encoder layers built a full
gene-by-gene attention matrix that only the one graph-regularized layer reads. Each was
replaced by a form that computes the same thing, held equal to the original by a test: the
operator and the decode on batched tensors, the ranks on the GPU, the fused attention kernel
on the unregularized layers, and the splits held in memory after one pass. After the changes
the stack samples show the process waiting on the GPU, on the operator's feed-forward block
over every gene of every strain; what is left is the architecture's own work.
\src{notes/experiments.019-simb-multimodal.scripts.gh_profile_cgt_expr.md}

\begin{table}[htbp]
  \centering
  \footnotesize
  \caption[]{\textbf{The same 1{,}200-epoch run is four times faster on Delta and the
  batch-128 form is twice as fast again.} Seconds per epoch per run, from
  \texttt{perf/epoch\_seconds} over at least 20 epochs, with the runs sharing a card as
  stated. \emph{Before} is the code v21 was launched on.}
  \label{tab:night-timing}
  \begin{tabular}{@{}llrrr@{}}
    \toprule
    machine and card & run & runs per card & s per epoch, before & s per epoch, after \\
    \midrule
    Delta, A40 & reference, batch 32 & 4 & 177 & 45 \\
    GilaHyper, RTX 6000 Ada & reference, batch 32 & 3 & 29 to 32 & 24 \\
    GilaHyper, RTX 6000 Ada & batch 128 & 3 & 19 to 20 & 11.4 \\
    \bottomrule
  \end{tabular}
\end{table}

A batch-128 run holds 12.2 GB of card memory against 8.5 GB at batch 32, so three fit on a
card and a fourth runs out of memory. Processes sharing a card take turns on it, so a card
of mixed arms runs every arm at the pace of the slowest; each card below holds one arm.

\pfstep{The state of v21 on Delta} The queue's start estimate for a resubmission was a day
out, so one pack was replaced inside its own running job at 03:58: the four reference runs
of split seeds 0 to 3 were stopped at epoch 70 and restarted from scratch on the new code
with zero loader workers and the splits in memory. They log 38 to 41 s per epoch, 100 s on
every tenth, about 45 s on average, which is 15 hours for 1{,}200 epochs. The other 26
packs were left as launched. Replacing them the same way is prepared and was run without
effect as a check; it needs the author's decision, as does the alternative of canceling and
resubmitting through the queue.
\src{experiments/019-simb-multimodal/scripts/delta_v21_inplace_swap.sh; W\&B runs iuqwl2p8,
tdnhpkid, fcq0xkke, 0jajrluv}

\pfstep{The rising validation loss on expression is a scale problem} Earlier sections
record that the validation loss bottoms out early and rises while the per-feature Pearson
keeps climbing, and that at its Pearson peak no run beats the per-gene mean on squared
error. A predictor with correlation $r$ to its target minimizes squared error when the ratio
of its spread to the target's equals $r$. On the six saved validation prediction sets, one
scalar was fit on half the validation strains (with a per-gene intercept) and scored on the
other half (\cref{tab:night-shrink}). The expression predictions are 1.4 to 2.2 times too
spread for their correlation; shrinking them by the fitted scalar brings all three
checkpoints below the per-gene mean, to what the correlation allows (for split seed 0,
$1.036 \times (1 - 0.225^2) = 0.983$ against a measured $0.984$). The proteome head has the
opposite sign at its best-validation checkpoint and is under-dispersed. So the ordering the
metric reads is intact and the magnitudes are overfit; calibration can be repaired after
training without changing the Pearson. The same arithmetic states how little is explained:
$r^2$ is 2 to 5 percent of per-gene variance across held-out strains.

\begin{table}[htbp]
  \centering
  \footnotesize
  \caption[]{\textbf{One scalar fit on half the validation strains brings every expression
  checkpoint below the per-gene mean on squared error.} NMSE is the squared error per gene
  over that gene's variance across the scoring strains, averaged over genes; 1 is the
  per-gene mean of the scoring strains, and the mean estimated on the fitting half reads
  above 1. Each row is a best-validation checkpoint (a selected epoch) of round v13 or v14;
  halves swapped and averaged. The lowest NMSE of each row is bold.}
  \label{tab:night-shrink}
  %% SOURCE: generated by experiments/019-simb-multimodal/scripts/v22_tables.py from results/prediction_shrinkage_probe.json
\begin{tabular}{@{}lrrrrrrr@{}}
\toprule
checkpoint & \shortstack[r]{val\\strains} & \shortstack[r]{val Pearson\\per feature} & \shortstack[r]{sd(pred) /\\sd(target)} & \shortstack[r]{fitted\\scalar $c$} & \shortstack[r]{NMSE,\\per-gene\\mean} & \shortstack[r]{NMSE,\\model\\as is} & \shortstack[r]{NMSE,\\shrunk\\by $c$} \\
\midrule
expression, V\_ref\_s2 & 155 & 0.144 & 0.317 & 0.46 & 1.030 & 1.052 & \textbf{1.016} \\
expression, V\_ref\_s0 & 155 & 0.225 & 0.319 & 0.70 & 1.036 & 0.993 & \textbf{0.984} \\
expression, V\_concat\_s0 & 155 & 0.192 & 0.343 & 0.52 & 1.036 & 1.029 & \textbf{0.997} \\
proteome, P\_concat\_s0 & 448 & 0.113 & 0.062 & 1.75 & 1.010 & \textbf{0.997} & 1.002 \\
proteome, P\_ref\_s0 & 448 & 0.116 & 0.085 & 1.38 & 1.010 & \textbf{0.995} & 0.997 \\
proteome, P\_ref\_s2 & 448 & 0.128 & 0.063 & 1.89 & 1.008 & 0.992 & \textbf{0.987} \\
\bottomrule
\end{tabular}

\end{table}

\pfstep{Round v22: what a faster batch costs} Every arm is the v21 reference with one
change, on the new code, 1{,}200 epochs unless stated, scored as v21 by the mean validation
Pearson per feature over epochs 1{,}000 to 1{,}200 (\cref{tab:night-arms,tab:night-runs}).
Two to three split seeds per arm size an effect and do not decide one.

\begin{table}[htbp]
  \centering
  \footnotesize
  \caption[]{\textbf{Batch 128 trains as far as the reference in half the time when its
  learning rate is left alone or warmed up, and collapses when it is scaled without
  warmup.} A run is \emph{collapsed} when its prediction spread ratio never reached 0.05 in
  300 epochs, or reached it and then stayed below 0.01 for 50 consecutive epochs, the
  per-gene-mean predictor; such runs were stopped. The highest arm mean is bold; an arrow
  marks an arm whose complete runs are all above or all below the reference on shared split
  seeds.}
  \label{tab:night-arms}
  %% SOURCE: generated by experiments/019-simb-multimodal/scripts/v22_tables.py from results/v22_readout.json
\begin{tabular}{@{}lp{46mm}rrrrrr@{}}
\toprule
arm & change against the reference & \shortstack[r]{split\\seeds\\run} & \shortstack[r]{collapsed\\or never\\launched} & \shortstack[r]{complete\\at 1{,}200\\epochs} & \shortstack[r]{mean val Pearson\\per feature, epochs\\1{,}000 to 1{,}200\\(complete runs)} & \shortstack[r]{mean paired\\difference vs F\_ref\\($n$ shared split seeds)} & \shortstack[r]{s per\\epoch} \\
\midrule
F\_ref & the v21 reference: batch 32, lr $3\times10^{-4}$, six layers, width 90 & 3 & 0 & 3 & 0.062 & +0.000 (3) & 24.0 \\
F\_b128lr1 & batch 128, lr $3\times10^{-4}$ & 3 & 0 & 3 & \textbf{0.075} & +0.013 (3) $\uparrow$ & 11.4 \\
F\_b128 & batch 128, lr $6\times10^{-4}$ & 3 & 2 & 1 & 0.059 & +0.000 (1) & 10.2 \\
F\_b128wu & batch 128, lr $6\times10^{-4}$ after a 50-epoch warmup & 3 & 0 & 3 & 0.067 & +0.005 (3) & 11.4 \\
F\_b128lr4 & batch 128, lr $1.2\times10^{-3}$ & 3 & 3 & 0 &  &  & 12.0 \\
F\_l4w180 & batch 32, four layers, width 180 & 2 & 2 & 0 &  &  & 17.6 \\
F\_b128wu\_hadam & F\_b128wu with the Hadamard operator & 3 & 0 & 3 & 0.036 & -0.025 (3) $\downarrow$ & 11.2 \\
F\_b128wu\_wd & F\_b128wu with weight decay 0.3, 800 epochs & 3 & 0 & 0 &  &  & 11.4 \\
F\_b128wu\_drop & F\_b128wu with dropout 0.3, 800 epochs & 3 & 1 & 0 &  &  & 11.4 \\
F\_b256lr1 & batch 256, lr $3\times10^{-4}$, 1{,}000 epochs & 1 & 0 & 0 &  &  & 3.1 \\
\bottomrule
\end{tabular}

\end{table}

\begin{table}[htbp]
  \centering
  \footnotesize
  \setlength{\tabcolsep}{3pt}
  \caption[]{\textbf{Every run of v22, with its W\&B link.} The three ladder columns are the
  mean validation Pearson per feature over the 20 epochs ending at that epoch, so runs
  stopped early can still be compared at matched epochs. A run with no entry in the window
  column did not finish 1{,}200 epochs; the state column gives the last epoch reached in
  parentheses. The highest window score is bold.}
  \label{tab:night-runs}
  %% SOURCE: generated by experiments/019-simb-multimodal/scripts/v22_tables.py from results/v22_readout.json
\begin{tabular}{@{}llrrrrrrrl@{}}
\toprule
 & & & \multicolumn{3}{c}{\shortstack{val Pearson per\\feature, mean of the 20\\epochs ending at epoch}} & & & & \\
\cmidrule(lr){4-6}
arm & W\&B run & \shortstack[r]{split\\seed} & 200 & 600 & 1{,}199 & \shortstack[r]{val Pearson,\\mean over\\epochs 1{,}000\\to 1{,}200} & \shortstack[r]{spread\\ratio, last\\epoch} & \shortstack[r]{train Pearson,\\last eval-mode\\pass} & \shortstack[l]{state (last\\epoch)} \\
\midrule
F\_ref & \href{https://wandb.ai/zhao-group/torchcell_019_expr_v22/runs/o053grx1}{\texttt{o053grx1}} & 0 & 0.004 & 0.007 &  &  & 0.437 & 0.512 & running (1{,}034) \\
F\_ref & \href{https://wandb.ai/zhao-group/torchcell_019_expr_v22/runs/7cmz3r5i}{\texttt{7cmz3r5i}} & 1 & -0.017 & 0.026 &  &  & 0.305 & 0.451 & running (1{,}029) \\
F\_ref & \href{https://wandb.ai/zhao-group/torchcell_019_expr_v22/runs/60puupp0}{\texttt{60puupp0}} & 2 & 0.095 & 0.033 &  &  & 0.347 & 0.443 & running (1{,}030) \\
\addlinespace
F\_b128lr1 & \href{https://wandb.ai/zhao-group/torchcell_019_expr_v22/runs/671rcqge}{\texttt{671rcqge}} & 0 & 0.014 & 0.013 & 0.076 & 0.069 & 0.401 & 0.525 & complete \\
F\_b128lr1 & \href{https://wandb.ai/zhao-group/torchcell_019_expr_v22/runs/jock61jk}{\texttt{jock61jk}} & 1 & -0.014 & 0.005 & 0.083 & 0.075 & 0.314 & 0.509 & complete \\
F\_b128lr1 & \href{https://wandb.ai/zhao-group/torchcell_019_expr_v22/runs/rza3mz2y}{\texttt{rza3mz2y}} & 2 & 0.115 & 0.032 & 0.085 & \textbf{0.081} & 0.444 & 0.535 & complete \\
\addlinespace
F\_b128 & \href{https://wandb.ai/zhao-group/torchcell_019_expr_v22/runs/l885wzyv}{\texttt{l885wzyv}} & 0 & 0.028 & 0.005 &  &  & 0.038 & 0.062 & collapsed at 212 (625) \\
F\_b128 & \href{https://wandb.ai/zhao-group/torchcell_019_expr_v22/runs/zz9sbrr3}{\texttt{zz9sbrr3}} & 1 & 0.006 & -0.006 &  &  & 0.000 & 0.019 & collapsed at 658 (930) \\
F\_b128 & \href{https://wandb.ai/zhao-group/torchcell_019_expr_v22/runs/z5fxcrnx}{\texttt{z5fxcrnx}} & 2 & 0.103 & 0.026 & 0.065 & 0.059 & 0.435 & 0.524 & complete \\
\addlinespace
F\_b128wu & \href{https://wandb.ai/zhao-group/torchcell_019_expr_v22/runs/49zyp3vw}{\texttt{49zyp3vw}} & 0 & 0.007 & 0.019 & 0.087 & 0.081 & 0.447 & 0.592 & complete \\
F\_b128wu & \href{https://wandb.ai/zhao-group/torchcell_019_expr_v22/runs/yokxkvjc}{\texttt{yokxkvjc}} & 1 & -0.019 & 0.019 & 0.085 & 0.071 & 0.358 & 0.527 & complete \\
F\_b128wu & \href{https://wandb.ai/zhao-group/torchcell_019_expr_v22/runs/i4n88yyy}{\texttt{i4n88yyy}} & 2 & 0.103 & 0.042 & 0.061 & 0.049 & 0.355 & 0.461 & complete \\
\addlinespace
F\_b128lr4 & \href{https://wandb.ai/zhao-group/torchcell_019_expr_v22/runs/rbn6guq7}{\texttt{rbn6guq7}} & 0 & 0.003 &  &  &  & 0.000 & 0.021 & never launched (442) \\
F\_b128lr4 & \href{https://wandb.ai/zhao-group/torchcell_019_expr_v22/runs/jzx4kljl}{\texttt{jzx4kljl}} & 1 & 0.002 &  &  &  & 0.000 & 0.007 & never launched (454) \\
F\_b128lr4 & \href{https://wandb.ai/zhao-group/torchcell_019_expr_v22/runs/a9mz22oq}{\texttt{a9mz22oq}} & 2 & 0.001 &  &  &  & 0.000 & 0.018 & never launched (446) \\
\addlinespace
F\_l4w180 & \href{https://wandb.ai/zhao-group/torchcell_019_expr_v22/runs/pbqtgnpe}{\texttt{pbqtgnpe}} & 0 & -0.004 &  &  &  & 0.000 & 0.008 & collapsed at 215 (312) \\
F\_l4w180 & \href{https://wandb.ai/zhao-group/torchcell_019_expr_v22/runs/ntn060cz}{\texttt{ntn060cz}} & 1 & -0.007 &  &  &  & 0.000 & -0.153 & collapsed at 119 (305) \\
\addlinespace
F\_b128wu\_hadam & \href{https://wandb.ai/zhao-group/torchcell_019_expr_v22/runs/56m6vbt8}{\texttt{56m6vbt8}} & 0 & 0.029 & 0.005 & 0.067 & 0.048 & 0.374 & 0.509 & complete \\
F\_b128wu\_hadam & \href{https://wandb.ai/zhao-group/torchcell_019_expr_v22/runs/j3l6huyu}{\texttt{j3l6huyu}} & 1 & -0.013 & -0.007 & 0.019 & 0.013 & 0.270 & 0.410 & complete \\
F\_b128wu\_hadam & \href{https://wandb.ai/zhao-group/torchcell_019_expr_v22/runs/wcp7wryw}{\texttt{wcp7wryw}} & 2 & 0.099 & 0.047 & 0.044 & 0.048 & 0.411 & 0.499 & complete \\
\addlinespace
F\_b128wu\_wd & \href{https://wandb.ai/zhao-group/torchcell_019_expr_v22/runs/18gw3p46}{\texttt{18gw3p46}} & 0 & 0.011 &  &  &  & 0.209 & 0.274 & running (443) \\
F\_b128wu\_wd & \href{https://wandb.ai/zhao-group/torchcell_019_expr_v22/runs/k38jsllm}{\texttt{k38jsllm}} & 1 & -0.020 &  &  &  & 0.190 & 0.264 & running (437) \\
F\_b128wu\_wd & \href{https://wandb.ai/zhao-group/torchcell_019_expr_v22/runs/t9avzqjw}{\texttt{t9avzqjw}} & 2 & 0.089 &  &  &  & 0.191 & 0.270 & running (430) \\
\addlinespace
F\_b128wu\_drop & \href{https://wandb.ai/zhao-group/torchcell_019_expr_v22/runs/w6sy4dlf}{\texttt{w6sy4dlf}} & 0 & 0.022 &  &  &  & 0.160 & 0.257 & running (441) \\
F\_b128wu\_drop & \href{https://wandb.ai/zhao-group/torchcell_019_expr_v22/runs/w08c6cda}{\texttt{w08c6cda}} & 1 & -0.013 &  &  &  & 0.000 & 0.108 & collapsed at 213 (436) \\
F\_b128wu\_drop & \href{https://wandb.ai/zhao-group/torchcell_019_expr_v22/runs/jo6tr18b}{\texttt{jo6tr18b}} & 2 & 0.110 &  &  &  & 0.082 & 0.224 & running (429) \\
\bottomrule
\end{tabular}

\end{table}

%%NIGHT-FINDINGS
